\begin{abstract}
We consider the problem of space-efficient result caching in search and recommender systems. A \emph{result cache} is a key-value data store mapping entities like queries or user ids to a list of associated items. With the heavy computational cost of modern retrieval pipelines, particularly in the age of large language models (LLMs), precomputing and caching results is often a necessity. However, result caching shifts the resource burden from compute to storage. Prior work on result caching has not addressed the distinct challenge that such caches must store a \emph{list} of values per key, rather than a single entity.
Our key insight is that result caches often exhibit strong power-law behavior, with frequent repetition of items across different keys. We propose to exploit this property by modeling the cache as a matrix — where rows correspond to keys and columns to ranked positions — and building \emph{compressed static functions} for each column. 

This architecture raises open research questions about codebook sharing across columns, entropy-reducing permutations for unordered results, and the handling of ragged lists of varying length.
In this paper, we systematically analyze each of these questions through a theoretical and experimental lens, culminating in the design of \sys, a new system for low-memory result caching in retrieval pipelines. By extending compressed static functions to handle lists of values, \sys performs lookups directly on highly compressed representations of the result cache while maintaining microsecond-level query latency. On workloads derived from real production data of a major e-commerce platform, \sys achieves an XXX reduction in memory compared to in-memory caching baselines, yielding a Pareto-optimal tradeoff between latency and memory footprint across the systems evaluated.
\end{abstract}